\documentclass[12pt,a4paper]{article}

\usepackage[utf8]{inputenc}
\usepackage[T1]{fontenc}
\usepackage[brazil]{babel}
\usepackage{amsmath}
\usepackage{geometry}

\geometry{
    top=2.5cm,
    bottom=2.5cm,
    left=3cm,
    right=3cm
}

\begin{document}

\subsection*{2.5. Criação da variável \texttt{low\_usage}}

\textbf{Com base nos quartis calculados anteriormente, utilize o primeiro quartil $Q_1$ como critério para classificar um dia como de ``baixa utilização'' do sistema. A partir desse critério, crie a variável binária \texttt{low\_usage}, determine o número e a proporção de dias classificados como de baixa utilização e explique o significado desse critério no contexto dos dados.}

Para realizar a classificação dos dias de acordo com o nível de utilização do sistema, foi utilizado como ponto de corte o primeiro quartil da variável \texttt{total\_user}.

O primeiro quartil, representado por $Q_1$, corresponde ao valor abaixo do qual se encontra aproximadamente 25\% das observações da distribuição. A partir dos dados analisados, obteve-se:

\[
Q_1 = 2105{,}5
\]

Dessa forma, foi adotado o seguinte critério para a construção da variável binária \texttt{low\_usage}:

\[
\texttt{low\_usage} =
\begin{cases}
1, & \text{se } \texttt{total\_user} < 2105{,}5,\\
0, & \text{caso contrário}.
\end{cases}
\]

Portanto, os dias em que o número total de usuários foi inferior a $2105{,}5$ foram classificados como dias de baixa utilização do sistema.

Como a variável \texttt{total\_user} representa uma contagem de usuários e assume valores inteiros, a condição

\[
\texttt{total\_user} < 2105{,}5
\]

equivale, na prática, a:

\[
\texttt{total\_user} \leq 2105.
\]

A variável \texttt{low\_usage} foi criada em R utilizando a seguinte expressão:

\begin{verbatim}
dados$low_usage <- ifelse(
  dados$total_user < Q1,
  1,
  0
)
\end{verbatim}

Após a aplicação desse critério às 300 observações presentes no conjunto de dados, foram encontrados:

\[
n_{\text{low\_usage}} = 75
\]

dias classificados como de baixa utilização.

A proporção desses dias em relação ao total de observações foi calculada por:

\[
p =
\frac{n_{\text{low\_usage}}}{n}
\]

Substituindo os valores encontrados:

\[
p =
\frac{75}{300}
\]

\[
p = 0{,}25
\]

Convertendo essa proporção para porcentagem:

\[
0{,}25 \times 100 = 25\%
\]

Assim, \textbf{75 dos 300 dias analisados foram classificados como dias de baixa utilização}, correspondendo a \textbf{25\% das observações}.

Esse resultado é coerente com a utilização do primeiro quartil como ponto de corte, uma vez que $Q_1$ delimita aproximadamente os 25\% menores valores da distribuição.

Nesse contexto, a variável \texttt{low\_usage} permite identificar de maneira objetiva os dias nos quais o sistema apresentou níveis relativamente baixos de utilização em comparação com o restante do período analisado.

Dessa forma, a interpretação da nova variável é:

\begin{itemize}
    \item \texttt{low\_usage = 1}: o dia apresentou baixa utilização, isto é, o valor de \texttt{total\_user} foi inferior a $2105{,}5$;
    
    \item \texttt{low\_usage = 0}: o dia não foi classificado como de baixa utilização, pois seu valor de \texttt{total\_user} foi igual ou superior a $2105{,}5$.
\end{itemize}

A nova variável foi acrescentada ao conjunto de dados original, preservando as demais variáveis existentes e permitindo que essa classificação seja utilizada nas análises posteriores.

\end{document}